In the strategy comparison, both proposed configurations achieved higher macro-$F_1$ scores than BLINKER and MNE on Raja, Cao2018, and the pooled evaluation (Table~\ref{tab:exp1_main}). Proposed-Med achieved macro-$F_1$ scores of 0.8825, 0.8068, and 0.8403, respectively, whereas Proposed-Mean achieved 0.8742, 0.8030, and 0.8345. The corresponding scores were 0.7654, 0.6940, and 0.7256 for BLINKER and 0.6664, 0.5627, and 0.6086 for MNE. This performance gap is reflected in the precision--recall profiles shown in Figure~\ref{fig:condition_prf}. BLINKER achieved the highest recall among all conditions on Raja, Cao2018, and the pooled evaluation, with values of 0.9518, 0.9715, and 0.9627, respectively, but also exhibited the lowest precision, at 0.6794, 0.5696, and 0.6182. Wilcoxon signed-rank tests on session-level $F_1$, with Bonferroni correction across the six pairwise comparisons, showed significant differences between each proposed configuration and each baseline. Proposed-Med versus BLINKER yielded $p=1.32\times10^{-11}$, and Proposed-Med versus MNE yielded $p=6.51\times10^{-9}$. Similarly, Proposed-Mean versus BLINKER yielded $p=4.97\times10^{-10}$, while Proposed-Mean versus MNE yielded $p=9.94\times10^{-9}$. In contrast, neither within-pair comparison was significant: BLINKER versus MNE yielded $p=0.812$, whereas Proposed-Mean versus Proposed-Med yielded $p=1$. Furthermore, no channel group in either dataset produced an inversion in which a baseline matched or exceeded the performance of Proposed-Med (Table~\ref{tab:exp2_inversions}).
